\section{Challenge Submissions}
\label{sec:challenge_submissions}

\subsection{Challenge subtask types}

The challenge is organized as a series of sets of tasks using
increasingly realistic representations of the data. In general, each
set of tasks includes 3 subtasks.

\begin{enumerate}
  \item{Estimate either per-object $p(z)$ or ensemble $n(z)$
      distributions for a set of different scenarios and provide the
      estimates in a specified format.}
  \item{Provide trained models for the different scenarios and a
      Python function that can be used to generate the estimates from subtask
      1 on an arbitrary dataset.}    
  \item{Provide a Python function that can be used to
      generate the models and estimates from subtasks 1 and 2 on
      arbitrary datasets.}
\end{enumerate}

The $p(z)$ estimates in subtask 1 and the trained models in subtask 2 should be provided
in a compressed \texttt{tar} file, which is described below.   Templates and instructions for the
Python functions needed for subtasks 2 and 3 will be provided and are
described below.


\subsection{\texorpdfstring{Data format for per-object $p(z)$ estimates}{Data format for per-object p(z) estimates}}

The $p(z)$ estimates should be submitted in \texttt{qp} format,
which allows users to specify a complete $p(z)$ distribution for
each object, as well as summary statistics for each object.

The \texttt{qp} package~\cite{QP} supports several different representations of
$p(z)$, such as different functional forms as well as interpolated
grids, histograms, and others.

For users unfamiliar with \texttt{qp}, we highly recommend representing
the $p(z)$ as either an interpolated grid or a Gaussian mixture model.

\begin{verbatim}
# Interpolated grid
import qp
import numpy as np
# Define the x-grid. Note that we put all the
# p(z) on the same x-grid
xvals = np.array([0,0.5,1,1.5,2])
# Define the y-values. Note we provide n_grid_points x n_objects 
# values, as we need to provide a y-value at each grid point 
# for each object.
yvals = np.array(
 [
   [0.01,0.2,0.3,0.2,0.01],
   [0.1,0.3,0.5,0.2,0.05]
 ]
)
ensemble = qp.interp.create_ensemble(xvals,yvals)
ensemble.write_to(<output_filename.hdf5>)
\end{verbatim}

\begin{verbatim}
# Mixture model
import qp
import numpy as np
# Define the means, standard deviations, and weights.
# These should each have shape n_objects, n_components.
# In this case we are defining 3 objects with 2-Gaussian 
# representations.
# For each object the weights should sum to 1, or they
# will be normalized.
means = [[0.3, 0.4], [0.5, 0.5], [0.6, 0.8]]
stds = [[0.2, 0.4], [0.1, 0.3], [0.05, 0.3]]
weights = [[0.8, 0.2], [0.7, 0.3], [0.8, 0.2]]
ensemble = qp.mixmod.create_ensemble(means=means,stds=stds,weights=weights)
ensemble.write_to(<output_filename.hdf5>)
\end{verbatim}

The submission files should use the same file name conventions defined
in Tab.~\ref{tab:file_fields}. The labels will typically be
\texttt{pz\_estimate} or \texttt{pz\_model} and will be specified in
the descriptions of the various tasks, e.g.,

\begin{verbatim}
\texttt{pz\_challenge\_taskset\_1\_cardinal\-\_pz\_estimate\_yr1.hdf5}
  or
\texttt{pz\_challenge\_taskset\_1\_cardinal\_pz\_model\_yr1.pkl}.
\end{verbatim}

All of these files should then be joined into a \texttt{tar} file, which 
should then be placed somewhere it can be download.  The URL for the
\texttt{tar} should be specified in
\texttt{tests/test\_\{submission\}.py}

\begin{verbatim}
SUBMISSION_NAME = "example"
SUBMISSION_URL = "https://your.institution.edu/submit_example.tgz"
\end{verbatim}


\subsection{Format for estimation-only Python functions and trained models}

For the second subtask, submissions should provide trained models and
implement a function to run estimation using those trained models on
the test files provided for each task set. The function will look something like this:

\begin{verbatim}
def run_taskset_1_estimation_only(
    model_file: str | Path,
    test_file: str | Path,
    output_file: str | Path,
) -> None:
    # do stuff and write p(z) estimates to "output_file"
\end{verbatim}

or

\begin{verbatim}
def run_taskset_2_estimation_only(
    model_file: str | Path,
    test_file: str | Path,
    output_file: str | Path,
) -> None:
    # do stuff and write p(z) estimates to "output_file"
\end{verbatim}

Templates for these functions are provided in the file
\texttt{tests/test\_\{submission\}.py} created as part of the setup.


\subsection{Format for training and estimation Python functions}

For the third subtask, submissions should implement a function to
train models and run estimation using those trained models on the
training and test files provided for each task set. The function will
look something like this:

\begin{verbatim}
def run_taskset_1_training_and_estimation(
    train_file: str | Path,
    test_file: str | Path,
    output_file: str | Path,
) -> None:
    # train a model using the "train_file" and make p(z) estimates 
    # and write them to "output_file"
\end{verbatim}

or

\begin{verbatim}
def run_taskset_2_training_and_estimation(
    train_file: str | Path,
    test_file: str | Path,
    output_file: str | Path,
) -> None:
    # train a model using the "train_file" and make p(z) estimates
    # and write them to "output_file"
\end{verbatim}

Templates for these functions are provided in the file
\texttt{tests/test\_\{submission\}.py} created as part of the setup.


\subsection{Submission mechanism}

Submissions will take the form of a pull request on the
\texttt{pz\_data\_challenge} repository and will include:

\begin{enumerate}
  \item{A file \texttt{tests/test\_\{submission\}.py} that includes the
      URL from which the compressed \texttt{tar} file should be downloaded
      as well as the Python functions for subtasks 2 and 3.   When
      created this will contain empty placeholder functions that will
      need to be implemented.}
   \item{A file \texttt{requirements\_\{submission\}.txt} that should
       be modified to include \texttt{pip} package names of any
       packages that need to be installed in order to run the
       functions in subtasks 2 and 3.}     
   \item{A file \texttt{.github/workflows/submit\_\{submission\}.yaml} to
       run the submission validation in a GitHub action.  This should
       not need to be modified unless the prerquisited installation
       requires more than just \texttt{pip}  installing packages.}
\end{enumerate}

All three of these files are created by the
\texttt{scripts/prepare\_submission.py} script.  

You will need modify the \texttt{tests/test\_\{submission\}.py} to
give the location of the \texttt{tar} file containing the PZ estimates
and trained models, and to implement the require functions.

See \texttt{https://github.com/LSSTDESC/pz\_data\_challenge/pull/6}
for an example of a submission.


\subsection{Submission validation}

The wrapping functions provided in the \texttt{tests/test\_\{submission\}.py}
file implement a number of \allowbreak checks on the data. Specifically, for each
expected file they check that:

\begin{enumerate}
    \item{the file exists;}
    \item{the file contains a valid \texttt{qp} ensemble;}
    \item{the \texttt{qp} ensemble includes ancillary data;}
    \item{the ancillary data includes a 'zmode' column with redshift
        estimates;}
    \item{the ancillary data includes an 'object\_id' column;}
    \item{the object\_ids in the submission file match the associated
        test file.}
\end{enumerate}

If any of these checks fail, the GitHub action triggered by the
submission will fail and report the cause of the failure.

The easiest way to test that you have correctly implemented the
required functions is simply to run these commands.

\begin{verbatim}
# Make sure that you have installed any packages you need
pip install -r requirement_{submission_name}.txt

# Run the functions you have provided as unit tests
py.test tests/test_{submission_name}.py
\end{verbatim}

if this succeeds, you can use a provided script to help you open the
pull request for your submisison.

\begin{verbatim}
# run the submission helper script.
python scripts/submit.py {submission_name}
\end{verbatim}

Note that the help script only prints the required commands, it does
not run them.  In short the command are:

\begin{verbatim}
# Check status of your local git clone by running git status, and make
# sure that you are on the branch submit/{submission_name} and do not
# have any files added or modified
git status

# Add your files to git
git add .github/workflows/submit_example.yaml
  requirements_example.txt
  tests/test_example.py

# Commit your files to your branch: 
git commit -m "Submitting {submission_name}"
  .github/workflows/submit_{submission_name}.yaml
  requirements_{submission_name}.txt
  tests/test_{submission_name}.py

# Push your commit
git push --set-upstream origin submit/{submission_name}

# Pushing to git should give you a URL that you can visit to create a
# pull request, for example:
#   https://github.com/LSSTDESC/pz_data_challenge/pull/new/submit/example
# Visit that URL and create a pull request, then add the 'submission'
# label to the PR.
# Finally, make sure that the github action validating your submission
# succeeds and fix any issues.
\end{verbatim}


\subsection{Submission aids}

A few scripts are provided to help you.

\begin{itemize}
  \item{\texttt{scripts/download\_public.py}:  downloads and unpacks
      the public data.}    
  \item{\texttt{scripts/prepare\_submission.py}:  sets up your area for
      a submission, creates the needed files from templates and
      downloads the public data, and suggests that you create a branch
      for you submission.}
  \item{\texttt{scripts/remove\_submission\_files.py}:  removes the
      submission files if you need to start over.}
  \item{\texttt{scripts/run\_metrics.py}:  run perfomance metrics on
      files in a submission you have created.}    
  \item{\texttt{py.test tests/test\_\{submission\_name\}.py}:  validates
      all the parts of your submission, checking that you have created
      all the required files and that they are properly formatted.}
\end{itemize}
      

  


      